\documentclass[11pt]{article}
\usepackage[a4paper,margin=2cm]{geometry}
\usepackage[hidelinks]{hyperref}
\usepackage{enumitem}
\setlist{itemsep=1pt,topsep=2pt}
\pagestyle{empty}
\begin{document}
\begin{center}
{\Large\bfseries LongChirp in one page}\\[2pt]
A test bench for finding very long, faint gravitational-wave chirps in real detector noise\\[2pt]
Sumit Pal, independent researcher (AI-assisted work) \quad Draft, 28 November 2026
\end{center}

\paragraph{The question.} If two black holes lighter than the Sun orbit each other and merge, they send out a faint
``chirp'' of gravitational waves. No star can make black holes that light, so finding one would point to something
new, such as black holes formed in the early universe. Their chirps last minutes in the LIGO detectors, not seconds,
and searching for them is expensive. Machine-learning searches are a cheaper idea, but they have mostly been tested on
short chirps: the main shared test so far used chirps of at most about 20 seconds.

\paragraph{What we built.} A fixed test, like an exam with an answer key:
\begin{itemize}
\item 27 hours of real LIGO data from 2019 (from the public Gravitational Wave Open Science Center);
\item 1,344 simulated chirps hidden in it at known times and distances;
\item 27 days of ``fake'' data, made by shifting the two detectors against each other, to measure how often noise
alone looks like a signal;
\item a scoring program: hand it a list of candidate detections and it says how far away a method can see, at a
fixed false-alarm rate (one false alarm per day, or per week).
\end{itemize}

\paragraph{What we found.} Three reference results, scored the same way:
\begin{itemize}
\item an ideal method that knows every hidden chirp's exact shape sees out to about 155 megaparsecs;
\item a standard template-matching search (3,235 templates) sees 57.5 megaparsecs at one false alarm per day;
\item a neural network sees 48.4 megaparsecs, and only 25.5 if it must say when each merger happened to within a
second.
\end{itemize}
So template matching beats the neural network by 12--19\% in distance, and both reach only about a third of what is
possible. That gap is the room a new method can try to win.

\paragraph{How it was checked.} The scoring program was tested on made-up answers where the right score is known.
The way chirps are added was checked against PyCBC, a standard tool of the field. Every downloaded data file was
checked against the numbers the data centre publishes. The checks found real problems along the way, which are
recorded rather than hidden: a waveform tool returns 5 seconds of data after each merger, which our first version
did not expect; one of our own tests failed once and was redone properly; and memory that a widely used package keeps
and never returns was reported to its developers.

\paragraph{What we do not claim.}
\begin{itemize}
\item We did not find any real signal, and we set no limit on how many such black holes exist.
\item Our template search is not the best the field can do; professional searches use millions of templates.
\item The results come from one stretch of 2019 data, two detectors, and one simple kind of signal (no spin).
\item Only one set of background data was used, so the error bars leave out some uncertainty.
\item The work was done with substantial help from AI models (Anthropic's Claude), which wrote most of the code under
the author's direction. It has been checked by automated tests and by a separate AI session, not by an outside
expert.
\end{itemize}

\paragraph{Where to look.} Paper: \texttt{longchirp\_preprint.pdf}. Code and data definition: DeepStrain release
\texttt{longchirp-v1.0.0}, \href{https://doi.org/10.5281/zenodo.24399603}{doi:10.5281/zenodo.24399603}.
Claim-by-claim checks: \texttt{REFEREE\_GUIDE.md}.
\end{document}
